\appendix
\clearpage
\refstepcounter{chapter}
\phantomsection
\addcontentsline{toc}{chapter}{付録A　設定とエビデンス索引}
\markboth{付録A　設定とエビデンス索引}{付録A　設定とエビデンス索引}
{\sffamily\bfseries\fontsize{22}{28}\selectfont 付録A\par}
\vspace{2mm}
{\sffamily\bfseries\fontsize{18}{24}\selectfont 設定とエビデンス索引\par}
\vspace{5mm}

\section{計算環境とデータの役割}

\begin{table}[H]
\centering
\caption{三つの計算環境の役割}
\begin{tabularx}{\textwidth}{>{\bfseries}p{34mm}p{42mm}X}
\toprule
環境 & 主な役割 & 本報告での利用\\
\midrule
データセンター & 各種ロボットデータの保存、編集、公開、データ編集センター運用 & 有効 ALOHA プロジェクト、軌跡、フレーム、時間、基盤機能を読み取り専用で確認\\
ALOHA 学習機 & 本番学習、履歴、実機運用モデル、収集ソフト、実機運用 & 学習済みモデル、学習履歴、アテンション記録、収集能力を読み取り専用で確認；ロボットは起動していない\\
ローカル研究機 & 強化学習、デジタルツイン、オフライン解析、報告生成 & 既存研究成果を読み取り、報告出力だけを新規作成\\
\bottomrule
\end{tabularx}
\end{table}

\section{基礎模倣データの項目}

\begin{table}[H]
\centering
\caption{基礎模倣学習データの主要項目}
\begin{tabularx}{\textwidth}{>{\bfseries}p{35mm}p{32mm}X}
\toprule
項目 & 形状・型 & 内容\\
\midrule
上部画像 & 映像フレーム & 作業台全体と双腕\\
左手首画像 & 映像フレーム & 左グリッパとボトル本体の近景\\
右手首画像 & 映像フレーム & 右グリッパとボトル口の近景\\
下部画像 & 映像フレーム & 原データには保存、本番運用モデルでは不使用\\
ロボット状態 & 14次元連続量 & 両腕関節と両グリッパ位置\\
動作 & 14次元連続量 & 次制御時刻の関節・グリッパ目標\\
タスク指示 & 文字列 & 当該軌跡の作業内容\\
時刻・フレーム番号 & 数値 & 順序、同期、軌跡長の検査\\
\bottomrule
\end{tabularx}
\end{table}

\section{主要エビデンス番号}

\begin{table}[H]
\centering\small
\caption{結論・エビデンス索引（E01--E08）}
\begin{tabularx}{\textwidth}{>{\bfseries}p{14mm}p{44mm}X}
\toprule
番号 & 種別 & 支持する内容\\
\midrule
E01 & 現場タスク説明と運用設定 & 左でボトル把持、右で開栓、本体とキャップを別箱へ投入\\
E02 & 収集ソフトの現行機能と履歴 & 連続収集、再収録、非同期保存、操作入力、健全性確認、停止\\
E03 & データセンターのプロジェクト・軌跡統計 & 51有効プロジェクト、2,413軌跡、2,907,804フレーム、約16.15時間\\
E04 & 確定版学習データの説明情報 & 25件の一意なデータ保管単位、1,051軌跡、879,852フレーム、項目整合\\
E05 & 全数値データ検査 & 非有限値、全要素がゼロの動作レコード、時刻、番号、完全重複の確認\\
E06 & 本番学習設定とモデル構成 & 3視野、状態・動作14次元、将来50制御周期分の動作列、学習目的\\
E07 & 現場運用設定と制御過程 & 50 Hz制御、先行再計画、動作制限、グリッパ保護\\
E08 & 元学習履歴 & 6,059記録、0--59,990ステップ、約51.4時間、損失曲線\\
\bottomrule
\end{tabularx}
\end{table}

\begin{table}[H]
\centering\small
\caption{結論・エビデンス索引（E09--E16）}
\begin{tabularx}{\textwidth}{>{\bfseries}p{14mm}p{44mm}X}
\toprule
番号 & 種別 & 支持する内容\\
\midrule
E09 & 運用モデルの保存情報 & ステップ19,000運用、8.38億パラメータ、正規化統計、読み取り可否\\
E10 & 現場観測声明 & 全工程、1時間超、約2本/分の概算と、映像・計数不足の明記\\
E11 & 条件・指示文の交差確認 & キャップなし、方向、反転の規模とキャップなし指示条件の不足\\
E12 & アテンション記録 & 9件の実行記録一覧、8,223サンプル、視点別構成比、解釈限界\\
E13 & 学習基盤の実行系譜 & 41試行、14設定、実行状態、到達ステップ\\
E14 & 実機設備写真 & 教示腕、作業腕、カメラ、作業領域；静止休止状態\\
E15 & 強化学習データ統計 & 洗浄835軌跡、切出しフレーム、238,816学習区間、人手評価\\
E16 & 強化学習保存モデルと履歴 & 6,154万パラメータ、ラウンド28--32、オフライン動作誤差\\
\bottomrule
\end{tabularx}
\end{table}

\begin{landscape}
\section{主要数値の再確認}
\begin{table}[H]
\centering\small
\caption{報告主要数値と集計範囲}
\begin{tabularx}{\linewidth}{>{\bfseries}p{66mm}p{48mm}X}
\toprule
指標 & 値 & 集計範囲・解釈\\
\midrule
データ基盤の有効 ALOHA データ & 2,413軌跡／16.15時間 & 51有効プロジェクト；プロジェクト外562軌跡を含めない\\
本番運用モデルの一意な学習データ & 1,051軌跡／4.89時間 & 25件の一意な確定版データ\\
フィルタ後の実学習データ & 844,102フレーム／4.69時間 & 50フレーム/秒換算\\
本番学習の等価サンプリング量 & 1,495軌跡相当 & 反復提示を含み、新規軌跡数ではない\\
実機運用モデル & 838,358,468パラメータ & モデルのみ；最適化状態なし\\
本番学習履歴 & 6,059記録／59,990ステップ & 約51.4時間；現場運用は19,000\\
学習試行 & 41試行／14設定 & 実行試行数であり成功実験数ではない\\
アテンション記録 & 8,223サンプル & 9件の実行記録一覧；結果ラベルなし\\
現場連続稼働 & 1時間超（概算） & 完全映像、計数表、標準試験手順なし\\
現場平均処理速度 & 約2本/分（概算） & 信頼区間は算出不能\\
強化学習データ & 835軌跡／238,816学習区間 & 洗浄研究用で、基礎分別データとは別\\
強化学習オフライン改善 & 約6.7\% & 固定検証の動作誤差；実機成功率ではない\\
\bottomrule
\end{tabularx}
\end{table}
\end{landscape}

\section{ソフトウェア環境}

\begin{table}[H]
\centering
\caption{監査で確認したソフトウェア環境}
\begin{tabularx}{\textwidth}{>{\bfseries}p{33mm}p{35mm}X}
\toprule
分類 & 版・状態 & 用途\\
\midrule
Python & 3.11.14 & ALOHA 学習機\\
JAX / JAXLIB & 0.5.3 / 0.5.3 & 基礎モデルの学習・推論\\
PyTorch & 2.7.1 & 一部データ・強化学習処理\\
\makecell[l]{Flax・Orbax\\Optax} & 0.10.2 / 0.11.13 / 0.2.4 & モデル、保存、最適化\\
\makecell[l]{OpenCV\\LeRobot} & 4.11.0.86 / 0.4.2 & 画像とロボットデータ\\
Docker & 29.0.2 & サービス実行環境\\
ROS / Isaac Sim & 当該指標は未記録 & 本報告は仮想環境を達成根拠にしていない\\
\bottomrule
\end{tabularx}
\end{table}

\begin{landscape}
\section{追加が必要な情報}
\begin{table}[H]
\centering\small
\caption{追加が必要な主要エビデンス}
\begin{tabularx}{\linewidth}{>{\bfseries}p{23mm}p{67mm}X}
\toprule
優先度 & 不足情報 & 補充方法\\
\midrule
P0 & 1本ごとの標準試験台帳と完全映像 & 固定場面、逐次ID、同期映像、二名確認、変更不能な集計表\\
P0 & タスク・工程の自動成否判定 & キャップ有無・離脱、本体・キャップの投入、異常状態を記録\\
P0 & キャップ有無とボトル口の向きの対照試験 & 同型・同位置で対にし、条件別誤り率を記録\\
P1 & 保存モデル間の統一比較 & 複数候補を凍結し、同一試験集合で評価\\
P1 & 映像の全数復号と意味品質 & 全数復号、キャップ有無・方向・工程ラベルを追加\\
P1 & 接触計測 & キャップ角、グリッパ電流、力または触覚を記録\\
P1 & 入力取得から制御出力までの遅延 & 推論、通信、動作列の切替、制御欠落の分布を記録\\
P1 & データ基盤・公開データの画面記録 & 審査用ブラウザ経路を整備後、ログイン後の表示範囲だけを撮影\\
P2 & 複数の乱数初期値と独立試験集合 & 3種類以上の初期値、学習に入れない試験場面を固定\\
\bottomrule
\end{tabularx}
\end{table}
\end{landscape}

\section{再生成と媒体上の制約}

報告資料一式には、機械可読統計、学習済みモデル情報、グラフ元データ、図生成手順、結論・エビデンス対応表、読み取り専用監査結果を保存した。ロボットを起動せず、遠隔データを変更せずに統計と図を再生成できる。手順は同梱の再生成ガイドに記載する。

写真、教示フレーム、アテンションマップ、概念図は種類を明示し、実験写真を人工生成していない。完全な展示映像がないため動画紹介ページは設けていない。データ編集サイトと公開データの画面は、審査時のブラウザ接続経路を安全に再確認してから追加する対象であり、本版では実在統計と実教示フレームを代替資料として用いた。
